% Folder 134-2, batch 06 — pages 101–120 (position in the folder; see below).
% Pass: Opus 5.5 (claude-opus-5-5), 2026-10-15 — first pass, unchecked against the pages by a human.
%
% Special rules for folder 134-2 (issues #44 and #57, the latter as decided by the
% user on 2026-10-15) applied: Pursuing Stacks text (printed in Maltsiniotis (ed.),
% A la poursuite des champs, vol. I, SMF 2022) is not transcribed, one \page{} and
% one \note{} per page; texts sent to him are transcribed in full; texts by others
% not addressed to him are summarised.
%
% Numbering: \page{} numbers the position in the folder (as batch 4 of folder 1
% does, and as the tiles and the facsimile do). Up to p. 114 it agrees with the
% archivists' pencil. From p. 115 on the pencil runs one ahead: p. 115 here is
% pencilled 116, p. 120 is pencilled 121. No sheet pencilled 115 is in the scan.
% Since his p. 12 (p. 114) breaks off mid-sentence (« the basic notion of a
% coherator ») and the letter's end is not on p. 115, the sheet pencilled 115 —
% presumably his p. 12', the end of the P.S. — seems to be missing from the
% digitisation (inference; it may also be a numbering slip of the archivists).
%
% Contents of the batch:
%  - pp. 101-114 (his pp. 5' to 12): end of the second typed copy of the letter
%    to Daniel Quillen (« Dear Daniel »), Pursuing Stacks §§ 4 (end) to 13:
%    §§ 5-12 (19/20/21.2.1983, « 21.2. » on p. 106), § 13 dated 22.2.1983 (p. 112),
%    signed « Alexander » (p. 114), P.S. « (23.2) » begun on p. 114. This copy
%    duplicates pp. 14-36; the rule kept is « imprimé, donc cité et non transcrit »:
%    one note per page naming the §§, his page number and where his handwritten
%    corrections fall (not transcribed in full).
%  - p. 115 (pencilled 116): R. Brown's typed note introducing his informal
%    translation of the letter to Breen of 17-19 July 1975. Others' text, not
%    addressed to him: summarised (one note).
%  - p. 116 (pencilled 117): blank; absent.
%  - pp. 117-120 (pencilled 118-121): R. Brown, « COMMENTS », typed, 4 pp.,
%    Bangor, November 1982 — comments on the letters to Breen sent to him.
%    Transcribed in full. Hands: « demander réf. : Breen » with a vertical stroke
%    and « Saavedra » (black ink, p. 117) read as his; « Memoir SMF », « Thesis,
%    U of Michigan » (p. 117) and the closing comment « 11 » (p. 120) are in a
%    lighter, rounder hand in grey pencil, NOT his black ink and not attributed
%    here (\note, not \add); « U.K. » (p. 120) in blue ink, not his.
%    FLAG for the parent: the brief took the line at the foot of p. 120 to be his;
%    on the page it looks like the same pencil hand as « Memoir SMF ».
%
% Letters in this batch: letter to D. Quillen, 19-22.2.1983 (pp. 101-114, his,
% printed: noted only). No letter to Breen. Third-party texts: Brown's note
% (p. 115, summarised) and Brown's Comments (pp. 117-120, sent to him, in full).
\documentclass[11pt,a4paper]{article}
\input{../preamble/grothendieck}

\folder{134-2}
\batch{6}
\pages{101}{120}
\foldertitle{[Chapitre I : Take off (pages 1 à 65) et table des matières provisoire] : tapuscrits annotés (19/02-22/02/1983), note manuscrite (s.d.), copies de lettre (1975, s.d.).}
\dating{1975-[1983]}
\watermark{Édition de démonstration}

\begin{document}

\section{Lettre à Daniel Quillen, 19--22.2.1983 (second exemplaire, suite) — \emph{Pursuing Stacks}}

\page{101}
\note{sa p. « 5' » : § 4 (suite). Second exemplaire de la lettre à Quillen, qui double les pp.~14 à 36 ; texte de \emph{Pursuing Stacks}, imprimé, donc cité et non transcrit (Maltsiniotis (éd.), \emph{À la poursuite des champs}, vol.~I, SMF 2022). Corrections manuscrites : « finite inverse limits » puis « lim » au-dessus de « fiber products » biffé, « iterated » ajouté quatre fois devant « fiber products », « here » au-dessus de « are » ; en marge gauche, une remarque en anglais : il abandonnera désormais le qualificatif « iterated ».}

\page{102}
\note{sa p. « 6 » : fin du § 4 et § 5 (numéro « 5. » ajouté en marge, sur une frappe « 6 »). Ajouts « This step » au-dessus de « which » biffé, « system of », « letters », « in this respect » ; en marge gauche, écrite en travers, une remarque en anglais sur l'ordre de l'opération $\ast$ (« heuristically clear », mais demandant une preuve).}

\page{103}
\note{sa p. « 6' » : § 5 (suite), début du § 6 (numéro « 7. » en marge, corrigé de la frappe). Ajouts « (iterated) », « will turn out to be », « of these », « inductive », « quite » ; en marge gauche, une note datée « (23.2) », appelée par un astérisque après « I believe it is », qui revient sur cette conviction (« But no longer … »).}

\page{104}
\note{sa p. « 7 » : § 6 (suite) — décomposition cellulaire des $D_n$ et première idée des sommes amalgamées « standard ». Les deux figures (triangle et cercle) sont dessinées à la main, suivies de « which are circles » ; « formulation » remplace un mot biffé ; retouches d'indices ($n$, $2n+1$, $0 \leq j \leq n-1$).}

\page{105}
\note{sa p. « 7' » : § 6 (suite et fin). Ajouts « i.e. numbered », « $(i \in I)$ », « sub- », « $f_{n(i),j}$ », « object » au-dessus de « subset », « $(2 \leq i \leq N)$ » ; la dernière phrase est complétée à la main (choix de source ou but selon les signes $\varepsilon$, $\varepsilon'$). Bas de page blanc.}

\page{106}
\note{sa p. « 8 » (numéro non visible), datée « 21.2. » : §§ 7 et 8 (numéros « 8. » et « 9. » en marge). Ajout « however » ; un astérisque après « all categories $C_n$ » renvoie à une note en marge gauche, en face d'un trait vertical, qui déclare l'assertion fausse et renvoie au P.S., p.~12.}

\page{107}
\note{sa p. « 8' » : § 8 (suite). Ajouts « the bunch of », « (binary) », « two », « finite », « a "homotopy" » au-dessus d'un mot biffé, « implicit », « that » ; deux notes en marge gauche, en face de traits verticaux, qualifient d'inexact ce qui est dit des sommes amalgamées (« unaccurate », « unaccurate, see above! »).}

\page{108}
\note{sa p. « 9 » : fin du § 8 (l'histoire des sept mathématiciens) et § 9 (numéro « 10. » en marge). Ajouts « his own », « which », « (in the sense of the Y-theory) », « $\infty$- », « n ».}

\page{109}
\note{sa p. « 9' » : § 9 (suite). Ajouts « i.e. », « with », « I believe », « two » au-dessus de « categories ».}

\page{110}
\note{sa p. « 10 » : § 10 (numéro « 11. » en marge). Ajouts « of $\mathrm{Hom}(D_n, |K|)$ », « by », « an "allowable" » au-dessus d'un mot biffé, « merely », « in the $D_i$ », « after all », « undesirable » au-dessus de « extra ».}

\page{111}
\note{sa p. « 10' » : § 11 (numéro « 12. » en marge), fin de la journée du 21.2 (« Good night »). Ajouts « $M$- », « stepwise », « gerbs » deux fois au-dessus de mots biffés, « should ».}

\page{112}
\note{sa p. « 11 », datée « 22.2.1983 » : § 12 (numéro « 13. » en marge ; c'est le § 13 de l'édition imprimée) — terminologie : $n$-stacks, $n$-Gr-stacks, « coherators » (« cohéreurs »). Peu de corrections manuscrites.}

\page{113}
\note{sa p. « 11' » : § 12 (suite). Ajouts « fiber products » à deux endroits sur « amalgamated sums » biffé, « primitive », « intrinsic », « even » ; en marge gauche, en face d'un trait vertical, « false, see P.S. » ; plus bas, en marge, « nor even automorphisms! ».}

\page{114}
\note{sa p. « 12 » : fin du § 12, formule finale « Very cordially yours », signée à la main « Alexander » ; P.S. daté « (23.2) » sur l'erreur concernant les produits fibrés dans les $C_n$, interrompu en bas de page. Ajouts « implies », « still in embryo », « later », « practically » ; en marge gauche, une remarque en anglais (« Maybe false, … », sur la sous-catégorie $C_0$). La suite du P.S. n'est pas dans ce lot.}

\section{R. Brown, note sur sa traduction de la lettre à Breen}

\page{115}
\note{note dactylographiée de R. Brown (Bangor) : la suite est une traduction informelle, par lui avec l'aide de L. Breen (Rennes) et de J.-L. Loday (Strasbourg), de la partie mathématique de la lettre manuscrite de Grothendieck à Breen datée des 17--19 juillet 1975 ; il remercie Grothendieck d'en avoir communiqué les copies, ainsi que des deux lettres précédentes. Texte d'un tiers, résumé. Aucune marque de Grothendieck.}

\section{R. Brown, « Comments » (Bangor, novembre 1982)}

\page{117}
\note{tapuscrit de R. Brown, adressé à Grothendieck, en quatre pages numérotées « -2- » à « -4- » à partir de la deuxième.}

\textbf{COMMENTS}

\begin{enumerate}
  \item[1.] The topic of duality has progressed somewhat since 1975, see for example

  L. Begueri \note{à la suite, au crayon gris, d'une main qui n'est pas son encre noire : « Memoir SMF »}

  M. Bester \note{à la suite, de la même main au crayon : « Thesis, U of Michigan »}

  (These references are due to L. Breen.)

  \marginal{demander réf. : Breen.}\note{à l'encre noire, à droite des deux références, en face d'un trait vertical}

  \item[2.] The work of Mme Sinh referred to on p.~2 is as follows:

  H.X. Sinh, \emph{Gr-catégories}, Thèse Doctorat d'État Université Paris VII (1975).

  See also

  N.\ill{}. \add{Saavedra,} Rivano, \emph{Catégories tanakiennes}, Springer Lecture Notes in Mathematics 265 (1972).\note{l'initiale dactylographiée après « N. » est surfrappée ; « Saavedra, » est écrit à l'encre noire au-dessus, avec un signe d'insertion}

  Related to these notions is that of $+$-groupoid (=coherent symmetric monoidal groupoid with strict associativity and strict zero object) whose theory is developed also in:

  M. Takeuchi, 'On Villamayor and Zelinsky's long exact sequence', Mem. Amer. Math. Soc. 33 (1981) No 249.

  \item[3.] The description of groups in (1-cat) given on p.~2--3 is in effect an equivalence between G-groupoids (= group objects in the category of groupoids) and crossed modules. This is described also in:

  R. Brown and C.B. Spencer, 'G-groupoids, crossed modules and the fundamental groupoid of a topological group', Proc. Kon. Akad. v. Weten. 79(1976), 296 -- 302.

  and is believed to have been known to Verdier in 1965.

  \item[4.] The equivalence mentioned in the last comment has been generalised to
\end{enumerate}

\page{118}
equivalences between six categories as set out in the diagram,

\begin{tikzcd}[column sep=small, row sep=small, nodes={font=\scriptsize}]
\text{poly-T-complexes} \arrow[r, leftrightarrow, "(f)"] \arrow[d, leftrightarrow, dashed, "(f)"'] & \text{cubical T-complexes} \arrow[r, leftrightarrow, "(c)"] & \omega\text{-groupoids} \arrow[d, leftrightarrow, "(d)"] \arrow[dl, leftrightarrow, "(a)"] \\
\text{simplicial T-complexes} \arrow[r, leftrightarrow, "(e)"'] & \text{crossed complexes} \arrow[r, leftrightarrow, "(d)"'] & \infty\text{-groupoids}
\end{tikzcd}

\note{la flèche (f) verticale est en pointillé sur la page.}

The categories are defined and the equivalences proved in:

\begin{enumerate}
  \item[a)] R. Brown and P. J. Higgins, 'The algebra of cubes', J. Pure Appl. Alg. 21(1981) 233 -- 260.
  \item[b)] R. Brown and P. J. Higgins, 'Colimit theorems for relative homotopy groups', \emph{ibid.} 22(1981) 11 -- 41.
  \item[c)] R. Brown and P. J. Higgins, 'The equivalence of $\omega$-groupoids and cubical T-complexes', Cah. Top. Géom. Diff. (3\textsuperscript{e} Coll. sur les catégories, dédié a Charles Ehresmann) 22 (1981), 349 -- 369.
  \item[d)] R. Brown and P. J. Higgins, 'The equivalence of $\infty$-groupoids and crossed complexes' \emph{ibid}, 370 -- 386.
  \item[e)] N. K. Ashley, \emph{Simplicial T-complexes}, Ph.D. thesis, University of Wales (1978),(to be published in an issue of \textbf{Esquisses Mathématiques}, Amiens).
  \item[f)] D. W. Jones, Ph.D. thesis in preparation.
\end{enumerate}

Applications of the equivalence (a) are given in reference (b) above.

\begin{enumerate}
  \item[5.] However the notion of $\infty$-groupoid referred to in (4) is not adequate for the purposes considered by Grothendieck since, roughly speaking, crossed complexes describe the homotopy types only of spaces $X$ which fibre over a $K(\pi,1)$ with fibre a product of Eilenberg-MacLane spaces.
  \item[6.] A notion of multiple groupoid called an \emph{n-cat-group} has been described
\end{enumerate}

\page{119}
in:

J.-L. Loday, 'Spaces with finitely many homotopy groups', J. Pure Appl. Alg. 24(1982) 179 -- 202.

The n-cat-groups (which may be described as n-fold categories in the category of groups) are adequate for the description of truncated homotopy types. This is Loday's main result.

\begin{enumerate}
  \item[7.] An analysis of the structure of non-abelian chain complex on a length two normalised chain complex of a simplicial group is given in:

  D. Conduché, 'Modules croisés généralisés de longueur 2' (in preparation).

  \item[8.] By following through some of the equivalences in (4) above, we can associate to an $\infty$-groupoid $G$ a simplicial T-complex $NG$, whose underlying simplicial set could be called the \emph{nerve} of $G$. A direct description of $NG$ is that $(NG)_n$ is the set of $\infty$-groupoid maps to $G$ from $\sigma\underset{\sim}{\Delta}^n$, the homotopy $\infty$-groupoid of the standard n-simplex $\Delta^n$ (i.e. $\sigma\underset{\sim}{\Delta}^n$ has associated crossed complex $\pi\underset{\sim}{\Delta}^n$, the homotopy crossed complex of $\underset{\sim}{\Delta}^n$).

  Now a definition of $\infty$-category is also given in (4)(d) above, but the corresponding functor nerve: $\infty$-categories $\longrightarrow$ simplicial sets is still unknown.

  \item[9.] An indication of the use of the ideas of (4) in non-abelian situations is given in:
\end{enumerate}

\begin{enumerate}
    \item[a)] R. Brown and P.J. Higgins, 'Crossed complexes and non-abelian extensions', Int. Conf. on Category Theory, Gummersbach (1981), Springer L.N.M. (to appear).
    \item[b)] R. Brown, 'Non-abelian cohomology and the homotopy classification of maps', Proc. Conf, \emph{Méthodes d'algèbre homotopique en topologie}, Marseilles -- Luminy (1982). Astérixes (to appear)
\end{enumerate}

\page{120}
The latter paper uses the equivalence between crossed complexes and simplicial T-complexes (which generalises the Dold-Kan theorem on the equivalence of chain complexes and simplicial abelian groups) to give a notion of nerve, and so classifying space $BC$ of a crossed complex $C$. The classical Eilenberg-MacLane bijection
\[
[X, K(A, n)] \cong H^n(X, A)
\]
for a CW-complex $X$, is then generalised to a bijection
\[
[X, BC] \cong H^0(X, C)
\]
However, as pointed out in (5), these results are still "too abelian", and much more work is needed to realise Grothendieck's programme.

\begin{enumerate}
  \item[10.] The theory of tubular neighbourhoods has been developed in papers by D.A. Cox.

  'Algebraic tubular neighbourhoods I, II', Math. Scand. 42 (1978), 211 -- 228, 229 -- 242.

  Another of his papers seems relevant:

  'Homotopy theory of simplicial schemes', Compositio Math. 39(1979) 263 -- 296.
\end{enumerate}

R. BROWN

School of Mathematics and Computer Science, University College of North Wales, Bangor, Gwynedd. U.K.\note{« U.K. » ajouté à l'encre bleue}

November, 1982

\note{au pied de la page, au crayon gris, de la même main que « Memoir SMF » (p.~117), et non de l'encre noire de Grothendieck : « 11 A van Kampen Theorem for 2-cat-groups has now been proved by Brown and Loday. »}

\end{document}
